\clearpage
\SourceLeaf{065}\pdfbookmark[1]{Abschnitt II.}{sectio-secunda}
\begin{center}
{\large ABSCHNITT II.}\par\medskip
\emph{Beziehungen für einen einzelnen Ort im Raum.}\par\bigskip
\hypertarget{native-065-article047}{}\pdfbookmark[2]{47.}{article-047}47.
\end{center}
Im ersten Abschnitt haben wir die Bewegung der Himmelskörper auf ihren Bahnen behandelt, ohne die Lage dieser Bahnen im Raum zu berücksichtigen. Um diese Lage zu bestimmen und so die Orte eines Himmelskörpers zu beliebigen anderen Punkten im Raum in Beziehung setzen zu können, benötigen wir offenbar sowohl die Lage der Bahnebene gegenüber einer bekannten Ebene (zum Beispiel der Ebene der Erdbahn, der \emph{Ekliptik}) als auch die Lage der Apsiden in jener Ebene. Das lässt sich am bequemsten mit sphärischer Trigonometrie behandeln. Wir denken uns eine Kugelfläche mit beliebigem Radius um die Sonne als Mittelpunkt. Auf dieser beschreibt jede durch die Sonne gehende Ebene einen Großkreis und jede von der Sonne aus gezogene Gerade einen Punkt. Zu Ebenen und Geraden, die nicht durch die Sonne selbst gehen, ziehen wir parallele Ebenen und Geraden durch die Sonne. Die ihnen entsprechenden Großkreise und Punkte auf der Kugelfläche stellen dann auch die ursprünglichen Ebenen und Geraden dar. Man kann sich auch eine Kugel mit einem, wie man sagt, unendlich großen Radius denken, auf der parallele Ebenen und Geraden in gleicher Weise dargestellt werden.

Fällt die Bahnebene also nicht mit der Ebene der Ekliptik zusammen, so schneiden sich die diesen Ebenen entsprechenden Großkreise (die wir ebenfalls einfach Bahn und Ekliptik nennen) in zwei Punkten, den \emph{Knoten}. Im einen Knoten geht der Körper, von der Sonne aus gesehen, aus dem südlichen Gebiet über die Ekliptik in das nördliche über. Im anderen kehrt er aus diesem in jenes zurück. Der erste heißt \emph{aufsteigender}, der zweite \emph{absteigender Knoten}. Die Lage der Knoten auf der Ekliptik geben wir durch ihren Abstand vom mittleren
\SourceLeaf{066}Frühlingspunkt (ihre \emph{Länge}) an, in der Reihenfolge der Tierkreiszeichen gezählt. In Fig.~1 sei $\AscendingNode$ der aufsteigende Knoten, $A\AscendingNode B$ ein Teil der Ekliptik und $C\AscendingNode D$ ein Teil der Bahn. Die Erde bewege sich von $A$ nach $B$ und der Himmelskörper von $C$ nach $D$. Offenbar kann der sphärische Winkel, den $\AscendingNode D$ mit $\AscendingNode B$ bildet, von $0$ bis $180^0$ anwachsen. Darüber hinaus wäre $\AscendingNode$ nicht mehr der aufsteigende Knoten. Diesen Winkel nennen wir die \emph{Neigung der Bahn} gegen die Ekliptik. Ist die Lage der Bahnebene durch die Länge des Knotens und die Neigung der Bahn bestimmt, so brauchen wir nur noch den Abstand des Perihels vom aufsteigenden Knoten. Diesen zählen wir in der Bewegungsrichtung selbst und setzen ihn daher negativ oder zwischen $180^0$ und $360^0$ an, wenn das Perihel südlich der Ekliptik liegt. Ferner seien die folgenden Ausdrücke festgehalten. Die Länge eines beliebigen Punktes auf dem Bahnkreis wird von dem Punkt aus gezählt, der auf der Bahn ebenso weit rückwärts vom aufsteigenden Knoten liegt wie der Frühlingspunkt auf der Ekliptik rückwärts von demselben Knoten liegt. Die \emph{Länge des Perihels} ist daher die Summe aus der Länge des Knotens und dem Abstand des Perihels vom Knoten. Die \emph{wahre Länge} des Körpers \emph{auf der Bahn} ist die Summe aus der wahren Anomalie und der Länge des Perihels. Schließlich heißt die Summe aus der mittleren Anomalie und der Länge des Perihels die \emph{mittlere Länge}. Dieser letzte Ausdruck kann offenbar nur bei elliptischen Bahnen verwendet werden.

